\section{Introduction}
\label{sec:intro}
The effective response of a heterogeneous material emerges from the deformation of its constituents and their interactions with the microstructure. At finite strains, this interaction can produce anisotropy and coupling between loading modes that are difficult to capture with a predetermined macroscopic constitutive law. Computational homogenization addresses this difficulty by resolving a representative volume element (RVE), in which an imposed deformation drives a microscopic boundary-value problem whose solution yields the effective stresses and constitutive derivatives~\cite{Miehe1999,kouznetsova2001approach,geers2010multi}. The effective law is thus computed rather than prescribed in advance.

This flexibility comes at a computational cost. Characterizing the material over a range of deformations already requires many nonlinear finite element solves; using the same RVE within a structural calculation multiplies those solves across macroscopic integration points and load increments~\cite{feyel2003multilevel,geers2010multi}. Much of this repeated work concerns closely related microscopic states. Surrogate modeling seeks to exploit that relationship by computing representative solutions offline and reusing their information in less expensive online evaluations. The goal is to reduce the cost of these evaluations without losing the constitutive behavior that made the RVE description useful.

Two complementary approximation routes are relevant, but they play different roles in this work. Non-intrusive constitutive surrogates learn the effective response from RVE input--output data and dispense with the microscopic solve online~\cite{le2015computational}. Intrusive projection-based reduced-order models (PROMs) instead compress the microscopic displacement field and retain projected equilibrium equations~\cite{benner2015survey,yvonnet2007,Hernandez2014}, with hyperreduction limiting the local contributions required to evaluate them~\cite{hernandez2017}. This work follows the first route and learns the effective stored energy while embedding the constitutive restrictions needed for its mechanical use. The intrusive route is retained as a secondary reference because it preserves an approximate microscopic state and the underlying constituent law.

The central contribution is a family of learnable paired directional features for polyconvex anisotropic energies. Their orientations and exponents adapt to the RVE response, while their paired construction makes the stress vanish exactly at the undeformed state without compromising polyconvexity. Coupling these features to either an input-convex neural network (ICNN) or an input-convex Kolmogorov--Arnold network (ICKAN) gives two realizations of the same constitutive design. PROM-based surrogates of the same RVE then provide secondary accuracy--cost references. Two structural examples follow, a tensile coupon and a confined compression test beyond the data, to show how the learned energies behave once they leave isolated material evaluations.

\subsection{Energy-based constitutive learning and polyconvexity}
Direct constitutive learning moves the expensive microscopic solves to the data-generation stage and replaces them online by a trained material model. Neural homogenization has demonstrated this separation~\cite{Aldakheel2023}, making systematic sampling of representative microscopic responses an essential part of the offline workflow~\cite{kalina2023}. What the surrogate must represent depends on the material class. History-dependent responses may require learned history-to-response maps~\cite{Liu2022} or internal descriptors such as maximum directional stretches~\cite{Ghaderi2020}, while interacting scales may motivate staggered data generation and training~\cite{AsadFarhat2026}. The present work, by contrast, considers history-independent hyperelasticity. This removes memory from the state but requires the stress to derive from a stored-energy function.

Learning that function, rather than fitting each stress component independently, makes the stress work conservative and gives a major-symmetric constitutive tangent by differentiation. This is a central motivation for the mechanics-informed energy formulation of As'ad et al.~\cite{asad2022}. Expressing the energy through objective deformation measures then removes dependence on a superposed rigid rotation. In this way, physics-augmented neural networks (PANNs) shift part of the burden from learning mechanical identities from data to satisfying them through the representation itself~\cite{klein2022,linden2023}. Prediction accuracy remains to be established, but potential structure and objectivity need not depend on how densely the training domain has been sampled.

Preserving a potential does not, by itself, control its curvature. At finite strains, that control must also be imposed in appropriate variables. Convexity in Green--Lagrange strain~\cite{asad2022} is attractive because stress and tangent follow directly from the learned energy. Its implication for curvature with respect to the deformation gradient nevertheless contains a stress-dependent term~\cite{GaoNeffRoventaThiel2017}, so it is not a global polyconvexity guarantee. Conversely, Xu et al.~\cite{Xu2021} parameterize an incremental tangent through a lower-triangular factorization to enforce positive semidefiniteness. This controls the tangent matrix but does not define the scalar potential required for integrability. The modeling task is therefore to preserve both the potential and the desired finite-strain curvature condition within one construction.

Polyconvexity provides a route to that combination. It represents the energy as a convex function of deformation minors, yielding a sufficient condition for rank-one convexity~\cite{ball1977}. For the two-dimensional model considered here, the independent arguments are the deformation gradient and its determinant. This framework also underpins existence results when the additional coercivity and admissibility hypotheses are met~\cite{Ciarlet1988}. Polyconvexity serves here as a sufficient constitutive restriction, not as a necessary property of every homogenized material or a guarantee of structural uniqueness. The remaining challenge is to keep enough freedom within these restrictions to describe the anisotropic RVE response.

\subsection{Directional features and representational expressiveness}
\label{sec:discussion}
For anisotropy, this challenge is especially consequential, because a flexible scalar core can distinguish only the deformation states that its inputs distinguish. Structural tensors encode material directions in objective scalar quantities and underpin established anisotropic polyconvex energies~\cite{SchroderNeff2003}. The directional description is itself a constitutive choice. Models of fiber-reinforced tissue, for example, must account for a distribution of orientations rather than one perfectly aligned family when that dispersion affects the mechanical response~\cite{Gasser2006}. A learned feature map must likewise expose the relevant directional distinctions while remaining compatible with objectivity and convexity in the deformation minors. It is therefore part of the constitutive model, not merely a preprocessing step.

A convex scalar core does not resolve this choice. An ICNN provides a convex function~\cite{amos2017}, but a sufficient composition with nonlinear features requires those features to be convex in the independent deformation minors and the core to be componentwise nondecreasing~\cite{klein2022}. This monotonicity imposes an order on feature space, since $\bm z(\F_a)\leq\bm z(\F_b)$ componentwise implies $H(\bm z(\F_a))\leq H(\bm z(\F_b))$. Increasing the network width cannot reverse that ordering. Reference corrections and growth terms, added later to complete the energy, do not remove this restriction on the core contribution. A poorly chosen feature family can therefore impose a representational obstruction before network capacity or optimization becomes limiting.

The numerical evidence of Klein et al.~\cite{klein2022} illustrates this distinction. For their cubic beam-lattice data, an exactly objective invariant-based polyconvex model fails to represent the shear response, whereas a model formulated in $\F$, $\cof\F$, and $\det\F$ gives a substantially better fit while approximating objectivity through data augmentation. The same invariant-based model performs very well for their analytically generated transversely isotropic data. In these results, accuracy depends on whether the admissible inputs expose the distinctions required by the material response, rather than on polyconvexity alone. This motivates enlarging the feature family before concluding that a convex core is inadequate.

The natural response is to make the directional description adaptive without abandoning the admissible composition. Generalized-structure-tensor models calibrate the structural tensors together with the energy model, allowing the data to identify aspects of the anisotropy~\cite{Kalina2024}. Recent triclinic invariant constructions and group symmetrization further broaden the polyconvex descriptions available for finite material-symmetry groups~\cite{klein2026}. Greater directional freedom, however, must be reconciled with the undeformed state. A learned anisotropic contribution need not produce zero reference stress, while cancelling that stress with an arbitrary strain correction can violate polyconvexity or material symmetry. The invariant-based normalization and analytic growth terms of Linden et al.~\cite{linden2023} demonstrate why these operations must belong to the admissible constitutive construction. The feature family must therefore do two jobs at once: distinguish the relevant directional response and control its reference derivatives so that exact stress normalization remains compatible with polyconvexity. This is the role of the paired features introduced here.

While the feature map determines what the model can distinguish, the scalar core determines how those distinctions are combined into an energy. Prior constructions distribute this flexibility in different ways. Constitutive artificial neural networks (CANNs) learn coefficients and exponents within an interpretable catalog of admissible energy terms~\cite{Linka2023}, whereas neural ordinary differential equations can learn monotone energy derivatives and recover convex energy contributions by integration~\cite{Tac2022}. In the present work, an ICNN combines the paired features through nonnegative weights and fixed convex activations. A Kolmogorov--Arnold network (KAN) instead places trainable univariate functions on its edges, a flexibility explored in constitutive KANs~\cite{Abdolazizi2025}. The feature-composition argument holds only if these edge functions are convex and nondecreasing, as in monotone input-convex KANs (ICKANs)~\cite{thakolkaran2025}. A smooth ICNN and an integrated-hat ICKAN are compared under the same paired features and analytic completion. The comparison isolates how the scalar dependence is represented, while the directional information and the mechanical guarantees stay the same.

\subsection{Projection-based reduced-order models}
The preceding approach learns the effective energy and thereby removes the microscopic solve from the online stage. Projection-based reduction makes a different compromise, retaining the RVE boundary-value problem but seeking its solution in a space learned from previous microscopic states. In an affine PROM, those states define a basis, commonly obtained through proper orthogonal decomposition (POD), and projected equilibrium determines the coordinates of each new displacement field. Reduced finite-strain homogenization has shown that this strategy can reuse information from earlier RVE solves without replacing the constituent law~\cite{yvonnet2007,Hernandez2014,Hernandez2020}. PROMs therefore provide a natural physics-based reference for the present study, since they preserve more of the microscopic problem than an effective-energy surrogate at the price of an online equilibrium solve.

The affordability of that solve depends first on the geometry of the solution family. A small number of loading parameters does not necessarily imply that the corresponding microscopic states lie close to a small global linear space. When their Kolmogorov $n$-width decays slowly, many basis vectors, and hence many equilibrium unknowns, are required to achieve adequate accuracy. Considerable effort has been devoted to mitigating this Kolmogorov barrier with richer reduced representations, from local linear spaces~\cite{Amsallem2012} to nonlinear manifolds built with autoencoders~\cite{Lee2020}, quadratic corrections~\cite{Barnett2022}, or learned closures in a POD coordinate space~\cite{barnett2023,sibuet2025discrete,aresdeparga2026nonlinear,Chmiel2024}. The last of these, PROM--ANN~\cite{barnett2023}, offers the compromise sought here, because it enriches the geometry of the reduced state without giving up either its POD description or projected equilibrium. Equilibrium determines only a small set of primary coordinates, while a regression model, here an artificial neural network (ANN), reconstructs the secondary coordinates that populate a much richer POD expansion~\cite{aresdeparga2026nonlinear}. The microscopic displacement can thereby draw on many modes while the online nonlinear system remains small, which motivates the primary--secondary construction adopted for the RVE in this work.

Reducing the number of unknowns does not make the reduced equations inexpensive, because assembling their residual and tangent may still require constitutive evaluations over the complete microscopic mesh at every iteration. While hardware acceleration performs this full-mesh work faster~\cite{Fritzen2016}, hyperreduction reduces how much of it must be performed, turning a PROM into a hyperreduced PROM (HPROM). Methods that act directly on the projected equations~\cite{Grimberg2021,AresDeParga2023}, instead of first reconstructing high-dimensional quantities as discrete empirical interpolation does~\cite{Chaturantabut2010}, include the linear-programming empirical quadrature procedure~\cite{patera2017lp,yano2019lp}, energy-conserving sampling and weighting (ECSW)~\cite{Farhat2014,farhat2015}, and the empirical cubature method (ECM)~\cite{hernandez2017,Hernandez2020}. Each replaces full assembly by a nonnegative weighted sum of selected element contributions while leaving the local constituent evaluations intact. This work follows ECM because previous developments extend its fixed rule, in which one support and one weight vector serve the whole training space, toward movable points~\cite{Hernandez2024}, subspace-adaptive weights~\cite{Bravo2024}, and finally the manifold-adaptive weights ECM (MAW--ECM)~\cite{hernandez2026dimensional}, which is the path needed to couple hyperreduction with PROM--ANN. In MAW--ECM, the weights become fields over the latent coordinates, so the same primary coordinates that locate the microscopic state on the PROM--ANN manifold also select the rule used to evaluate its equilibrium and stress.

Together, these choices yield three intrusive references for the structural comparison. The affine HPROM retains a linear POD state, solves projected equilibrium, and evaluates it with fixed ECM. HPROM--ANN replaces the linear state by the primary--secondary manifold and the fixed rule by MAW--ECM, but still determines its primary coordinates through equilibrium. Its direct variant, D-HPROM--ANN, removes that final online solve by evaluating the decoder at the imposed strain, as in the hyperelastic example of~\cite{hernandez2026dimensional}. Moving across the three models removes successive portions of the online microscopic work while preserving an approximate microscopic state and selected evaluations of the constituent law. They contextualize the accuracy and cost of the proposed energy surrogates but do not share their constitutive guarantees, because the approximate effective stress of a reduced model need not derive from a single effective potential, even when the constituent law is retained and the cubature weights are nonnegative.

\subsection{Contribution and outline}
Against these intrusive physics-based references, the central route pursued in this work is an effective-energy surrogate whose anisotropy is learned without relinquishing polyconvexity or exact reference-state normalization. The construction rests on paired directional features with trainable orientations, stretch powers, and determinant exponents constrained by explicit joint-convexity bounds.

The pairing is designed so that the strain derivative of every feature at the undeformed state is proportional to the identity tensor, even though its response away from that state remains anisotropic. The complete learned contribution can consequently be normalized by a single scalar correction affine in the determinant. This correction cancels the reference stress without tying that cancellation to the independently adjustable volumetric stiffness, allowing the directional description to adapt to the RVE while the reference state remains under analytic control.

With the directional information in this form, two monotone convex cores provide alternative representations of its scalar energetic dependence: a smooth ICNN and an integrated-hat spline ICKAN. In both cases, the learned core, reference correction, and growth terms form one complete energy for which the analysis establishes objectivity, an exactly stress-free reference, two-dimensional polyconvexity, and the specified growth limits. Global nonnegative energy is treated separately through a sufficient lower bound evaluated for the trained parameters. The novelty lies in bringing established convex-network and affine-determinant normalization principles together through the paired anisotropic family and its explicit analytical conditions.

The assessment works at two levels. At the constitutive level, comparing the constrained energies on independent RVE responses with the Unconstrained energy, a network trained under the same protocol without the curvature and growth constraints, measures the accuracy that the guarantees cost. The same study compares learned with prescribed paired directions and locates where the Unconstrained energy loses rank-one convexity beyond the data, while checks of the tangent and of the sign of the energy probe behavior that a fit metric cannot reveal~\cite{Tac2024}. At the structural level, the learned energies replace the RVE query of the first-order two-level finite element formulation, denoted FE$^2$, in which an RVE problem supplies the material response at each macroscopic integration point~\cite{feyel2003multilevel}. In a tensile coupon, they are compared with a nested full-order model (FOM)--FE$^2$ reference, and the three reduced models described above serve as physics-based accuracy--cost benchmarks. A confined compression test then drives the same energies beyond the data to ask which of them keep the structural problem solvable.

The conclusions are tied to the two periodic cells, plane-strain kinematics, and the two structural loadings. Agreement with FOM--FE$^2$ measures surrogate error within the adopted homogenization model, not error relative to a structure in which every microscopic cell is resolved explicitly, and the reported timings cover the online workflows only.

The remainder of the paper is organized as follows. Section~\ref{sec:fom} defines the homogenization problem and work-conjugate variables. Section~\ref{sec:learned} introduces the admissibility criteria, constructs the learned energies, and proves their properties, after which Section~\ref{sec:reduced} defines the reduced microscopic references used in the structural comparison. Section~\ref{sec:cell} reports the numerical examples, first the multicavity RVE (MC-RVE) representation study and then the single-cavity RVE (SC-RVE) deployment study, which ends with the coupon FE$^2$ comparison and a confined compression test beyond the data. Section~\ref{sec:conclusions} closes the paper, and the supplementary material documents additional curvature and closed-cycle diagnostics.

\FloatBarrier
